% TOP_PROBLEM: {"id": 3312, "title": "Large induced forest in a planar graph.", "category_id": 3, "category": {"id": 3, "name": "graph_theory", "display_name": "Graph Theory", "description": "Problems involving graphs, networks, and their properties.", "slug": "graph-theory", "order_index": 3, "created_at": "2026-07-31T15:26:25.671Z"}, "status": "open", "problem_number": "OPG-46634", "set_id": 12}
% FIRST_POSTED: 2026-10-02
% ATTEMPT_STATUS: unresolved
% UnsolvedMath ID 3312; source catalogue code OPG-46634
% !TeX program = lualatex
% GPT-6 Astra Ultra research attempt, 2026-10-02; independently checked partial work.
% Completion estimate: no numerical percentage assigned; see final status and literature audit.
\documentclass[11pt,a4paper]{article}
\usepackage[margin=25mm]{geometry}
\usepackage{fontspec}
\setmainfont{lmroman10-regular.otf}[BoldFont=lmroman10-bold.otf,
 ItalicFont=lmroman10-italic.otf,BoldItalicFont=lmroman10-bolditalic.otf]
\usepackage{amsmath,amssymb,mathtools}
\usepackage{unicode-math}
\setmathfont{latinmodern-math.otf}
\AtBeginDocument{\let\setminus\smallsetminus}
\usepackage{xurl,hyperref}
\hypersetup{colorlinks=true,urlcolor=blue}
\setcounter{secnumdepth}{0}
\setlength{\parindent}{0pt}
\setlength{\parskip}{5pt}
\setlength{\emergencystretch}{3em}
\begin{document}
\section*{Large induced forest in a planar graph.}\label{3312}
{\small UnsolvedMath ID 3312 · OPG-46634 · Graph Theory · OpenGarden}
\subsection{Definitions and mathematical statement}
Let $G=(V,E)$ be a finite simple planar graph on $n$ vertices. For
$U\subseteq V$, its induced subgraph $G[U]$ has vertex set $U$ and
all edges of $G$ whose ends lie in $U$. A forest is an undirected
graph with no cycle; it may have several components and isolated
vertices. Define
\[
                      a(G)=\max\{|U|:G[U]\text{ is a forest}\}.
\]
The Albertson--Berman conjecture in this record is that
\[
                                a(G)\ge\frac n2
\]
for every finite simple planar graph. In integer terms this requires
at least $\lceil n/2\rceil$ vertices.

Equivalently, one should be able to delete a set $S\subseteq V$ of
at most $\lfloor n/2\rfloor$ vertices meeting every cycle. Such a set
is a feedback vertex set, and $G-S$ must retain every edge between
the remaining vertices. Selecting a spanning forest by simply
removing edges is therefore insufficient: a forest subgraph on all
vertices need not be induced.

The conclusion asks for one large induced forest, without requiring
that the complementary vertices also induce a forest. No girth,
maximum-degree, or connectedness assumption accompanies planarity.
The factor one half cannot be increased uniformly: a disjoint union
of copies of $K_4$ is planar and contributes at most two vertices
per copy to an induced forest, since every three vertices of a
$K_4$ already form a cycle. All bounds concern the number of
vertices rather than edges.
\subsection{Short English statement}
Can at least half of the vertices of every planar graph be retained while the induced graph has no cycle?
\subsection{Research attempt}
\paragraph{Scope and process.}
GPT-6 Astra Ultra, 2 October 2026. The canonical formulation above is retained. This attempt uses a primary-source literature audit, independent restricted proofs, and adversarial checks. Reproved results are not claimed as new. Numerical tests below are reproducible in \texttt{scripts/check\_attempt\_batch03\_extremal\_2026\_10\_02.py}.

\paragraph{Literature correction and scope.}
The catalogue records the Albertson--Berman half-order conjecture. It must not be described as currently open without qualification: Cames van Batenburg, Goedgebeur and Jooken [S3], posted 24 August 2026, report that it has been disproved, determine minimum counterexample order 29, and construct an infinite family with maximum induced-forest order at most $25n/52$. Their complete counterexample computation is not reproduced here. This attempt independently proves a substantial restricted case and identifies a necessary structural feature of any counterexample. The unresolved attempt status refers to our own complete verification, not to the historical conjecture's literature status.

\paragraph{A degeneracy criterion proved constructively.}
A graph is $d$-degenerate if every nonempty subgraph has a vertex of degree at most $d$. We prove that every $(2s-1)$-degenerate graph has a partition into $s$ induced forests. Delete vertices one by one with current degree at most $2s-1$, then insert them back in reverse order. At the moment a vertex $v$ is inserted, it has at most $2s-1$ already inserted neighbours. If each of the $s$ classes contained at least two of these neighbours, there would be at least $2s$ such neighbours. Hence some class contains at most one.

Assign $v$ to that class. Its induced subgraph was a forest by induction. Adding a vertex with at most one neighbour in that forest cannot create a cycle: a cycle through the new vertex would require two incident edges within the class. Other classes are unchanged. Induction proves the partition, and its largest class has at least $\lceil n/s\rceil$ vertices. In particular every 3-degenerate graph has an induced forest on at least $\lceil n/2\rceil$ vertices. This is a statement about induced subgraphs, since all edges between vertices of a chosen class are retained throughout.

\paragraph{Applying the criterion to planar graphs.}
Every triangle-free simple planar graph is 3-degenerate. To see this, take any nonempty subgraph. If it has at most two vertices, the degree conclusion is immediate. Otherwise the planar triangle-free edge bound $m\le2n-4$ gives average degree less than four, so some vertex has degree at most three. For disconnected subgraphs the same conclusion follows componentwise; the familiar edge bound can also be proved by Euler's formula after separately treating tree components. Thus the desired half-order statement holds for every triangle-free planar graph.

For arbitrary simple planar graphs, Euler's edge bound instead ensures 5-degeneracy, and the same argument gives a partition into three induced forests and an induced forest of order at least $\lceil n/3\rceil$. This is only an elementary lower bound, not a claim to the best bound in the literature. The difference between at most three and at most five previously inserted neighbours is exactly where the two-class proof ceases to work.

\paragraph{Structural and sharpness checks.}
If a graph violates the half-order bound, it cannot be 3-degenerate. It therefore has a nonempty subgraph of minimum degree at least four; equivalently its 4-core is nonempty. This is a necessary condition, not a sufficient recipe for a counterexample. For example, the octahedral graph has minimum degree four but an induced three-vertex path and no larger induced forest, giving equality at six vertices. Disjoint copies of $K_4$ also attain exactly one half: three vertices of a block already make a triangle, while two are always permitted.

A proper four-colouring is insufficient by itself. The union of two independent colour classes is bipartite, but bipartite induced subgraphs can contain cycles, as $K_{2,2}$ demonstrates. Thus grouping colour classes does not replace the degeneracy argument. The script checks the constructive two-forest partition on every 3-degenerate labelled graph of order at most five and verifies the stated extremal small examples by exhaustive induced-subset tests.

\paragraph{Final status.}
The half-order conclusion is independently established for all 3-degenerate graphs, including triangle-free planar graphs. The unrestricted historical claim is reported disproved by [S3]; this attempt has not independently reconstructed that full disproof. Its own full-target verification is therefore \textbf{UNRESOLVED}, with the precise gap being verification of a cited planar counterexample rather than a missing proof of a true universal theorem.

\subsection{Sources}
\begin{itemize}
\item[S1] ULAM.AI and the UnsolvedMath Contributors, supplied problems.json, record 3312 (OPG-46634), OpenGarden collection. Catalogue identity and source transcription; CC BY 4.0.
\url{https://huggingface.co/datasets/ulamai/UnsolvedMath}.
\item[S2] Open Problem Garden, Large induced forest in a planar graph.. Original problem and discussion; the mathematical formulation below is expanded from the supplied source record.
\url{http://www.openproblemgarden.org/op/large_induced_forest_in_a_planar_graph}.
\item[S3] W. Cames van Batenburg, J. Goedgebeur and J. Jooken, Counterexamples to the Albertson--Berman conjecture: minimum order, connectivity and an improved ratio bound, arXiv:2608.23260v1, 24 August 2026. Literature disproof report, not a counterexample certified by this attempt.
\url{https://arxiv.org/abs/2608.23260}.
\end{itemize}
\end{document}
